\lysh{15826_SOLUTION}{1}
{\bfseries ΛΥΣΗ}

{\bfseries Α)}
Η (1) είναι της μορφής $x^2 + y^2 + Ax + By + \Gamma = 0$, όπου $A = -2(\lambda + 1)$, $B = -2\lambda$ και $\Gamma = 2\lambda + 1$. Είναι:
\[ A^2 + B^2 - 4\Gamma = (-2(\lambda + 1))^2 + (-2\lambda)^2 - 4(2\lambda + 1) \]
\[ {} = 4\lambda^2 + 8\lambda + 4 + 4\lambda^2 - 8\lambda - 4 = 8\lambda^2 \]
Για να παριστάνει η (1) κύκλο πρέπει και αρκεί $A^2 + B^2 - 4\Gamma > 0 \Leftrightarrow 8\lambda^2 > 0 \Leftrightarrow \lambda \neq 0$.
Το κέντρο είναι το $K\left( - \frac{A}{2}, - \frac{B}{2} \right)$ δηλαδή $K(\lambda + 1, \lambda)$ και η ακτίνα
\[ \rho = \frac{\sqrt{A^2 + B^2 - 4\Gamma}}{2} = \frac{\sqrt{8\lambda^2}}{2} = \frac{2\sqrt{2}|\lambda|}{2} = \sqrt{2}|\lambda| \]

{\bfseries Β)}
Για $\lambda = 0$ η (1) γίνεται $x^2 + y^2 - 2x + 1 = 0 \Leftrightarrow (x - 1)^2 + y^2 = 0 \Leftrightarrow \left\{ \begin{array}{l} x = 1 \\ y = 0 \end{array} \right\}$
που σημαίνει ότι παριστάνει το σημείο $M(1, 0)$.

{\bfseries Γ)}
i. Η ευθεία που διέρχεται από τα κέντρα $K_1, K_2$ του σχήματος έχει συντελεστή διεύθυνσης $\lambda = \frac{y_{K_1} - y_{K_2}}{x_{K_1} - x_{K_2}} = \frac{2 - 1}{3 - 2} = 1$ και εξίσωση $\varepsilon: y - 1 = 1(x - 2) \Leftrightarrow y = x - 1$.

Θα αποδείξουμε ότι τα κέντρα όλων των κύκλων που προκύπτουν από την (1) βρίσκονται πάνω στην ευθεία $\varepsilon$. Πράγματι το τυχαίο κέντρο $K(\lambda + 1, \lambda)$ ανήκει στην ευθεία $\varepsilon$, αφού οι συντεταγμένες του επαληθεύουν την εξίσωση $y = x - 1$.
$\lambda = (\lambda + 1) - 1 \Leftrightarrow \lambda = \lambda$.

ii. Οι κύκλοι του σχήματος διέρχονται από το σημείο $M(1,0)$. Θα αποδείξουμε ότι όλοι οι κύκλοι που προκύπτουν από την (1) διέρχονται από το $M(1,0)$. Πράγματι οι συντεταγμένες του $M$ επαληθεύουν την (1) για κάθε $\lambda \in \mathbb{R}$, αφού
$1^2 + 0^2 - 2(\lambda + 1) \cdot 1 - 2\lambda \cdot 0 + 2\lambda + 1 = 0 \Leftrightarrow 1 - 2\lambda - 2 + 2\lambda + 1 = 0 \Leftrightarrow 0 = 0$.

iii. Θα πρέπει το κέντρο $K(\lambda + 1, \lambda)$ να απέχει από την ευθεία $\varepsilon: x + y - 1 = 0$ απόσταση ίση με την ακτίνα $\rho$.
\[ d(K, \varepsilon) = \frac{|\lambda + 1 + \lambda - 1|}{\sqrt{1^2 + 1^2}} = \frac{|2\lambda|}{\sqrt{2}} = \frac{2|\lambda|}{\sqrt{2}} = \sqrt{2}|\lambda| = \rho \]

Σημείωση: Η ευθεία $\varepsilon: x + y - 1 = 0$ διέρχεται από το σημείο $M(1,0)$ και να είναι κάθετη στην ευθεία $\zeta$, όπως φαίνεται και στο παρακάτω σχήμα.

% TODO: TikZ σχήμα
\end{lysh}